% Nouveau EGA V (cote 25): archive cote 25, site page 122 (PDF page 123).
% typescript. Draft compared: none.
% AI-transcribed by Claude Opus 5.5 (claude-opus-5-5, adaptive thinking, effort high), Message Batches API, batch ega25b (msgbatch_01A7n5YuXubUPNK1zsepi8RG), 2026-11-13; checked against both supplied scan parts; see CHECK_DECISIONS.tsv.
\scanpage{25/122}
\mspage{25}

systêmes linéaires de diviseurs sans composante fixe  sur X paramétrés
par $\overset{\hand{\vee}}{P}$ , et l'ensemble $\hand{\unsure{\psi|\emptyset}}$  des applications rationelles de X dans P, \hand{qui}
\inserted{sur $\bar{k}$ ne se factorisent par aucun hyperplan de $\unsure{\overline{P}|P}_{\bar{k}}$.}


En effet, \inserted[typed]{lorsque les anneaux locaux de X sont factoriels.} \deleted{\xout{X étant régulier, \unsure{il|X} en est de même de \unsure{$X\times_S P$} qui est lisse}
\xout{sur X, donc d'après un résultat de Auslander-Buchsbaum (qui sera repris}
\xout{au Chap III) \illegible{1} de les anneaux locaux de \unsure{$X\times_S P$} sont factoriels,}
et comme codim(X-U,X) \quad 2} , il s'ensuit que tout Module inversible sur
U peut se prolonger en un Module inversible sur X, donc la condition
envisagée dans b) est automatiquement satisfaite. D'autre part, d'après
\deleted{\unsure{un Mémoire} \illegible{1} \unsure{dans un} n° \unsure{ultérieur}}
Auslander-Buchsbaum (\deleted{qui figurera un jour au Chap III}\inserted{cf §20}) un annea\hand{/}u local
régulier est factoriel.

\margin{\unsure{$\times\quad\times\quad\nearrow$}}La proposition 2  b)  nous montre que myennant une hypothèse
\inserted[typed]{(savoir $D\quad\hand{\unsure{\phi|\emptyset}}$)}
supplémentaire sur un\xout{e famille} système linéaire de diviseurs D, on peut
lui associer canoniquement un Module inversible $\underline{L}$ sur X, prolongement
\inserted[typed]{où U=X-Z.}
canonique de $\underline{L}^{x}_{U} = f\ (\underline{O}_P(1))$ ,  On peut en donner une autre interprétation,
indépendante de l'hypothèse  $D\quad\hand{\unsure{\phi|\emptyset}}$ , de la façon suivante; par constructio

\xout{l'image inverse de $\underline{L}$x} ou presque $\underline{L}^{x}_{U}\otimes_S\underline{O}_P(1)$ \xout{n'}est \xout{autre que} \inserted[typed]{can isomorphe à} la restric-%
tion $\underline{N}'$ à $U\mathrm{x}_S\overset{\hand{\vee}}{P}$  du  Module inversible $\underline{N}$  sur $X\mathrm{x}_S\overset{\hand{\vee}}{P}$  défini par D. En vertu
du \xout{n°1} n°1, cela détermine \xout{uniquement $\underline{N}$, et} $\underline{L}_U$  à isom unique près;
et  (comme U est dense dans X)  il existe  un Module inversible $\underline{L}$ sur X,
déterminé à isom unique près (do\unsure{n|ç}c prolongeant $\underline{L}_U$ ) tel que  l'on ait
un isomorphisme  $\underline{L}\otimes_{\underline{O}_S}\underline{O}_{\overset{\hand{\unsure{\vee|/}}}{P}}\ (1) \quad \underline{N}$ ;  $\underline{L}$  est donné par exemple par
\[ \underline{L} \quad pr_1\ (\underline{N}) \quad ; \]
on l'appelle  le Module inversible sur X associé au systême lin de divi-%
suurs D.

\mspage{122}
